% Einheit 3 von 3 – Faktorisieren (bank/binomische-formeln/e3.jsonl, zusammenbau v0.1)
%% TODO Zweigzeile Teil 1: was hier gelernt wird (Satz in Schülersprache aus dem Einheitstitel) – Einheit 3
\einheitenkopf[e3]{Einheit 3 von 3 · Faktorisieren}
\zweigzeile{ab Kl. 8, je nach Buch bis Kl. 10 · keine P10-Aufgabe}
\setcounter{aufgabe}{16}

%% TODO Titel: Ich-kann-Satz fehlt in der Bank (Platzhalter: Kettenname) – Nr. 17
%% TODO Anweisung: ein Satz über der ersten Teilaufgabe fehlt in der Bank – Nr. 17
\begin{aufgabe}{Faktorisieren}
%% TODO \rechenplatz{n}: Schrittzahl des Grundfalls fehlt in der Bank (nur bei fester Schreibform, 2.3 b)
\begin{teile}
\teil $x^2 + 16x + 64$ – Quadrate mit Wurzeln? Passt das Mittelglied? \janein Wurzeln \leerfeld und \leerfeld
\end{teile}
\begin{gleichungsraster}[2]
\gl{\text{Faktorisiere: $x^2 - 36$}} &
\gl{\text{Faktorisiere: $x^2 + 8x + 16$}} \\
\gl{\text{Faktorisiere: $x^2 - 18x + 81$}} &
\gl{\text{Faktorisiere: $16x^2 - 9$}} \\
\gl{\text{Faktorisiere: $2x^2 - 32$}} &
\end{gleichungsraster}
\begin{teile}
\teil Lässt sich $x^2 + 64$ mit einer binomischen Formel faktorisieren? Begründe.
\teil Stimmt $x^2 - 16x + 64 = (x - 8)^2$? Mach die Probe durch Ausmultiplizieren.
\end{teile}
\begin{gleichungsraster}[2]
\gl{\text{Löse durch Faktorisieren: $x^2 - 144 = 0$}} &
\gl{\text{Schreib mit quadratischer Ergänzung als Quadrat plus Zahl: $x^2 + 8x + 3$}} \\
\gl{\text{Faktorisiere vollständig und mach die Probe: $3x^2 - 48$}} & \\
\end{gleichungsraster}
\end{aufgabe}

%% TODO Titel: Ich-kann-Satz fehlt in der Bank (Platzhalter: Kettenname) – Nr. 18
%% TODO Anweisung: ein Satz über der ersten Teilaufgabe fehlt in der Bank – Nr. 18
\begin{aufgabe}{Faktorisieren – Fehler finden, Begründen}
\begin{teile}
\teil Finde den Fehler und rechne richtig: \rechnung{x^2 + 1 &= (x + 1)^2}
\teil Warum lässt sich $x^2 + 16$ nicht mit einer binomischen Formel faktorisieren? Begründe.
\end{teile}
\end{aufgabe}
